\textbf{结论名称：} 弦切角定理\\
\textbf{适用条件：} 圆中，一直线与圆相切于弦的一个端点，该直线与弦所夹的角（弦切角），所对的弧为该弦所截得的某一弧。\\
\textbf{变量说明：} 设$\odot O$中，弦$AB$，直线$l$切圆于点$A$。在$l$上取点$T$（$T\neq A$），则$\angle BAT$为弦切角。点$C$为圆上不与$A,B$在$l$同一侧的任意一点，对应圆周角$\angle BCA$所对的弧为$\wideparen{AB}$（不含点$A$侧切线方向的弧）。\\
\textbf{核心公式：} \[ \boxed{\angle BAT = \angle BCA} \]\\
\textbf{使用场景：} 利用切线求圆周角，或反之判定切线；与圆幂定理结合证明角度相等、相似三角形；在切线与割线夹角计算中简化过程。\\
\textbf{简单例子：} 已知$PA$切$\odot O$于$A$，$AB$为圆的一条弦，点$C$在优弧$AB$上满足$\angle ACB = 55^\circ$，则由弦切角定理可直接得弦切角$\angle PAB = 55^\circ$。若测量得$\angle PAB = 55^\circ$，也可反推$PA$为切线。